% Ad Atticum 6.1 --- headnote
% ad-atticum-06-01, 24 February 50 BC, Laodicea

\textit{Cicero to Atticus, written from Laodicea on the
sixth day before the Kalends of March (24 February)
50 BC (the manuscript dateline: \textit{Scr.
Laodiceae vi K. Mart. a. 704 (50)}). The opening
letter of \textit{Ad Atticum} book 6, and the longest
and most substantial of Cicero's surviving despatches
from his proconsular year in Cilicia. He has come down
from the late-summer campaign on the Amanus to winter
in Laodicea, where the regional assizes occupy him
through the spring; the date of the letter is fixed in
its first sentence to five days before the Terminalia.
The letter answers, in order, an unusually full
bulletin from Atticus, and discharges the major
accumulated business of the governorship at one
sitting.}

\textit{The substantive heart of the letter --- sections
five through seven --- is Cicero's first detailed
handling of the Salaminian-loan affair, which becomes
a leitmotif of his correspondence for the rest of the
proconsulship. M. Iunius Brutus, behind the proxies
M. Scaptius and P. Matinius, had lent money to the
city of Salamis on Cyprus at four percent per month
compounded --- forty-eight percent a year --- under a
senatorial dispensation procured to evade the Lex
Gabinia's twelve-percent ceiling. The predecessor
Appius Pulcher had given Scaptius a prefecture and
squadrons of cavalry to enforce collection, with which
he had besieged the Salaminian senate in their own
council-house until five of them starved to death.
Cicero, holding himself to twelve percent uniformly
across the province by his own edict, has refused to
endorse the higher rate, refused the prefecture to
Scaptius as a man of business (per a standing
exception Atticus had himself negotiated), and
withdrawn the squadrons. He is now writing to Atticus
to set out the moral and legal accounting of the
choice --- and to brace him for the certainty that
Brutus, whose moneylending Cicero had not understood
to be personal until the affair began, will be
displeased. The tone is dry, organised, almost legal,
broken by flashes of indignation at Scaptius and of
weary irony at his own situation as the trustee of an
honour he had not realised he was being asked to
underwrite.}

\textit{Around this core the letter ranges across the
whole field of Cicero's current preoccupations: the
contrast between his administration and Appius's
(section two), the financial straits of the king
Ariobarzanes and the disproportionate share Pompey is
extracting from him (section three), the praetextate
education of the two Cicero boys at Laodicea (section
twelve), the impending Parthian invasion and the
auxiliaries of Deiotarus (section fourteen), the
two-part structure of his own provincial edict
(also section fourteen), his handling of the publicans
(section sixteen), the antiquarian quibbles raised by
Atticus about the statue of Africanus and the calendar
of Cn. Flavius (sections seventeen and eighteen), the
domestic question of Tullia's marriage (section ten),
the news of Vedius's baggage and the five portraits of
married ladies it contained (section twenty-five),
and a last whimsical thought about building a propylaeum
for the Academy at Athens to match Appius's at Eleusis
(section twenty-six). The political backdrop is the
deferral of the great question --- the consular
provinces, the Parthian command, the standing of
Caesar --- to the Kalends of March, the day after this
letter is written, with Cicero acutely afraid that his
own command will be extended and his return to Italy
prevented. The Greek-tag scaffolding is dense and
intimate throughout, in the inside-Atticus voice;
the closing date-formula --- ``the seven hundred and
sixty-fifth day after the battle of Leuctra'' --- is
the kind of antiquarian flourish the two of them
exchange for amusement.}
